\documentclass[11pt,oneside]{article}

% Standalone ProveIt research note.  The source intentionally has no external
% inputs so that the TeX/PDF pair is directly reproducible after extraction.
\usepackage[a4paper,margin=26mm,headheight=15pt]{geometry}
\usepackage[T1]{fontenc}
\usepackage[utf8]{inputenc}
\IfFileExists{libertinus.sty}{\usepackage{libertinus}}{\usepackage{lmodern}}
\usepackage{microtype}
\usepackage{amsmath,amssymb,amsthm,mathtools,mathrsfs,bm}
\usepackage{booktabs,longtable,tabularx,array}
\usepackage{graphicx}
\usepackage{xcolor}
\usepackage{enumitem}
\usepackage{fancyhdr}
\usepackage{xurl}
\usepackage{listings}
\usepackage[most]{tcolorbox}
\usepackage{hyperref}
\usepackage[nameinlink,noabbrev,capitalise]{cleveref}

\definecolor{linkblue}{RGB}{24,72,124}
\definecolor{deepgreen}{RGB}{23,102,69}
\definecolor{deepred}{RGB}{140,42,42}
\definecolor{deepgold}{RGB}{122,91,19}
\definecolor{shadegray}{RGB}{246,247,249}
\definecolor{rulegray}{RGB}{195,201,208}
\definecolor{palered}{RGB}{251,236,236}
\definecolor{palegold}{RGB}{251,245,230}
\definecolor{palegreen}{RGB}{233,245,243}
\definecolor{paleblue}{RGB}{236,243,251}

\hypersetup{
  hypertexnames=false,
  colorlinks=true,
  linkcolor=linkblue,
  citecolor=deepgreen,
  urlcolor=linkblue,
  pdftitle={Asymptotic Reversion of x + W(x) and a Calculus for Logarithmic Transseries},
  pdfauthor={Vladimir Reshetnikov},
  pdfsubject={Asymptotic expansion of the inverse of x plus Lambert W and general logarithmic-transseries reversion},
  pdfkeywords={Lambert W, generalized Lambert W, r-Lambert W, transseries, logarithmic asymptotics, series reversion, Lagrange-Burmann inversion}
}

\pagestyle{fancy}
\fancyhf{}
\fancyhead[L]{\small Reversion of \(x+W(x)\)}
\fancyhead[R]{\small\nouppercase{\leftmark}}
\fancyfoot[C]{\thepage}
\renewcommand{\headrulewidth}{0.4pt}

\setcounter{tocdepth}{2}
\setcounter{secnumdepth}{3}
\setlist[itemize]{leftmargin=1.5em,itemsep=0.25ex,topsep=0.6ex}
\setlist[enumerate]{leftmargin=1.9em,itemsep=0.35ex,topsep=0.6ex}
\allowdisplaybreaks[2]

\theoremstyle{plain}
\newtheorem{theorem}{Theorem}[section]
\newtheorem{proposition}[theorem]{Proposition}
\newtheorem{lemma}[theorem]{Lemma}
\newtheorem{corollary}[theorem]{Corollary}
\theoremstyle{definition}
\newtheorem{definition}[theorem]{Definition}
\newtheorem{algorithm}[theorem]{Algorithm}
\newtheorem{example}[theorem]{Example}
\theoremstyle{remark}
\newtheorem{remark}[theorem]{Remark}

\crefname{algorithm}{algorithm}{algorithms}
\Crefname{algorithm}{Algorithm}{Algorithms}

\newtcolorbox{mainbox}[1][]{%
  enhanced,breakable,colback=paleblue,colframe=linkblue,
  boxrule=0.7pt,arc=1.2mm,left=2mm,right=2mm,top=1.4mm,bottom=1.4mm,
  title={Main result},fonttitle=\bfseries,#1}
\newtcolorbox{warningbox}[1][]{%
  enhanced,breakable,colback=palered,colframe=deepred,
  boxrule=0.6pt,arc=1.2mm,left=2mm,right=2mm,top=1.2mm,bottom=1.2mm,
  title={Scale warning},fonttitle=\bfseries,#1}
\newtcolorbox{methodbox}[1][]{%
  enhanced,breakable,colback=palegreen,colframe=deepgreen,
  boxrule=0.6pt,arc=1.2mm,left=2mm,right=2mm,top=1.2mm,bottom=1.2mm,
  title={Reusable method},fonttitle=\bfseries,#1}
\newtcolorbox{summarybox}[1][]{%
  enhanced,breakable,colback=palegold,colframe=deepgold,
  boxrule=0.6pt,arc=1.2mm,left=2mm,right=2mm,top=1.2mm,bottom=1.2mm,
  title={Summary},fonttitle=\bfseries,#1}

\lstdefinestyle{pythonstyle}{%
  language=Python,
  basicstyle=\ttfamily\small,
  columns=fullflexible,
  frame=single,
  rulecolor=\color{rulegray},
  backgroundcolor=\color{shadegray},
  showstringspaces=false,
  keepspaces=true,
  breaklines=true,
  xleftmargin=0.4em,xrightmargin=0.4em,
  aboveskip=0.8em,belowskip=0.8em
}
\lstset{style=pythonstyle}

\newcommand{\W}{\operatorname{W}}
\newcommand{\rW}{\mathcal{W}}
\newcommand{\e}{\mathrm{e}}
\newcommand{\dd}{\mathop{}\!\mathrm{d}}
\newcommand{\R}{\mathbb{R}}
\newcommand{\C}{\mathbb{C}}
\newcommand{\Q}{\mathbb{Q}}
\newcommand{\Acal}{\mathcal{A}}
\newcommand{\Ccal}{\mathcal{C}}
\newcommand{\Tcal}{\mathcal{T}}
\newcommand{\ordt}{\nu_t}
\newcommand{\coeff}[2]{\left[#1^{#2}\right]}
\newcommand{\fallop}[2]{\left(#1\right)_{\!\underline{#2}}}
\DeclareMathOperator{\Res}{Res}
\DeclareMathOperator{\Trunc}{Trunc}

\begin{document}

\title{\textbf{Asymptotic Reversion of \(x+\W(x)\)}\\[0.35em]
  \Large A General Calculus for Logarithmic Transseries}
\author{Vladimir Reshetnikov\\
  \small ProveIt research note}
\date{2 September 2026}
\maketitle

\begin{abstract}
Let \(\W=\W_0\) be the principal real Lambert function and let
\(g=f^{-1}\), where \(f(x)=x+\W(x)\) on \([0,\infty)\).  This article derives
an all-orders asymptotic expansion of \(g(z)\) as \(z\to+\infty\), and uses the
example to develop a general method for reversing logarithmic transseries near
the identity.

The exact substitution \(u=\W(g(z))\) reduces the inverse problem to
\(u\e^u+u=z\).  Thus \(u\) is the \(r\)-Lambert function at \(r=1\), and
\(g(z)=z-u\).  If \(v=\W(z)\), then the useful asymptotic expansion is not
merely a series in inverse powers of \(\log z\), but the finer transseries,
which converges for all sufficiently large \(z\),
\[
 g(z)=z-v+\sum_{m\ge1}\frac{A_m(v)}{z^m},
\]
where every \(A_m\) is an explicit rational function.  We obtain the operator
formula
\[
 A_m(v)=\frac{(-1)^{m+1}}{(m+1)!}
 \prod_{j=0}^{m-1}\left(\frac{v}{1+v}\frac{\dd}{\dd v}-j\right)v^{m+1},
\]
a triangular recurrence, a direct Lagrange coefficient formula, and a
uniform remainder estimate of order \((\W(z)/z)^{N+1}\) after \(N\) correction
blocks.  A second reversion, in the parameters
\(q=\e^{-v}=v/z\) and \(s=1/v\), yields harmonic-number formulas for the
large-\(v\) structure of the coefficient blocks.

At the coarser logarithmic level, \(g(z)\) and \(z-\W(z)\) have identical
complete expansions in powers of \(1/\log z\); their difference is beyond all
orders of that scale.  We give the resulting Stirling-number formula and show
precisely how the exponentially small correction must be restored.  Finally,
for a general near-identity logarithmic transseries \(F(X)=X+A(X)\), we prove
the locally finite reversion identity
\[
 F^{-1}(X)=X+\sum_{n\ge1}\frac{(-1)^n}{n!}
 \frac{\dd^{n-1}}{\dd X^{n-1}}A(X)^n,
\]
derive an explicit block formula over an arbitrary differential coefficient
algebra, compare fixed-point, triangular, Lagrange--B\"urmann, and Newton
algorithms, and give residual certificates that turn formal truncations into
rigorous analytic approximations.
\end{abstract}

\noindent\textbf{Keywords.}
Lambert \(W\); generalized Lambert function; \(r\)-Lambert function;
logarithmic transseries; compositional inverse; Lagrange--B\"urmann inversion;
series reversion; asymptotic scale; harmonic numbers.

{\small\tableofcontents}
\clearpage

\section{The inverse problem and the answer in four forms}
\label{sec:intro}

Throughout, \(\W(x)=\W_0(x)\) is the principal real branch, characterized on
\([0,\infty)\) by
\begin{equation}
 \W(x)\e^{\W(x)}=x,\qquad \W(x)\ge0.
 \label{eq:Wdef}
\end{equation}
Define
\begin{equation}
 f(x)=x+\W(x),\qquad x\ge0,
 \label{eq:fdef}
\end{equation}
and write \(g=f^{-1}\).

The standard differentiation formula
\begin{equation}
 \W'(x)=\frac{\W(x)}{x(1+\W(x))}
       =\frac{\e^{-\W(x)}}{1+\W(x)},
 \qquad x>0,
 \label{eq:Wprime}
\end{equation}
with \(\W'(0)=1\), shows that \(f'(x)>0\).  Since \(f(0)=0\) and
\(f(x)\to\infty\), the inverse exists uniquely on \([0,\infty)\).

\begin{mainbox}
Put
\[
 v=\W(z),\qquad L=\log z,
 \qquad \ell=\log L,
 \qquad q=\e^{-v}=\frac{v}{z}.
\]
Then the inverse \(g=f^{-1}\) admits the following complementary descriptions.
\begin{enumerate}
\item \emph{Exact generalized-Lambert form:}
\[
 g(z)=z-\rW_1(z),
 \qquad \rW_1(z)\e^{\rW_1(z)}+\rW_1(z)=z.
\]
Here \(\rW_r\) denotes the \(r\)-Lambert function, not the complex branch
\(\W_1\) of the ordinary Lambert function.

\item \emph{Derivative or Lagrange--B\"urmann form:}
\begin{equation}
 g(z)\sim z+\sum_{n\ge1}\frac{(-1)^n}{n!}
 \left(\frac{\dd}{\dd z}\right)^{n-1}\W(z)^n.
 \label{eq:derivative-main}
\end{equation}
For this particular problem the formal expansion reorganizes into a genuinely
convergent local inverse series for all sufficiently large positive \(z\).

\item \emph{\(\W\)-anchored outer expansion:}
\begin{align}
 g(z)={}&z-v+\frac{v^2}{z(1+v)}
 +\frac{v^3(v^2-2)}{2z^2(1+v)^3}\notag\\
 &+\frac{v^4(2v-1)(v^3-6v-6)}{6z^3(1+v)^5}
 +O\!\left(\left(\frac{v}{z}\right)^4\right).
 \label{eq:anchored-first}
\end{align}
An all-orders coefficient formula and six explicit blocks are given below.

\item \emph{Pure logarithmic expansion:}
\begin{equation}
 g(z)\sim z-L+\ell-\sum_{n\ge1}\frac{P_n(\ell)}{L^n},
 \label{eq:pure-log-main}
\end{equation}
where
\begin{equation}
 P_n(y)=\sum_{k=1}^{n}(-1)^{n-k}
 \genfrac{[}{]}{0pt}{}{n}{n-k+1}\frac{y^k}{k!}
 \label{eq:Pn-stirling}
\end{equation}
and \(\genfrac{[}{]}{0pt}{}{n}{k}\) is an unsigned Stirling cycle number.
This last expansion is complete in the scale \(1/L\), but it omits the
beyond-all-orders difference between \(g(z)\) and \(z-\W(z)\).
\end{enumerate}
\end{mainbox}

The rest of the article proves these assertions, explains how they fit into a
single transseries, and extracts a reusable reversion calculus.

\section{Exact reduction to the \texorpdfstring{\(r\)-Lambert}{r-Lambert} function}
\label{sec:exact}

\subsection{The generalized Lambert function}

For a fixed real parameter \(r\), the \(r\)-Lambert function is locally defined
as the inverse of
\begin{equation}
 h_r(u)=u\e^u+ru.
 \label{eq:rLambert-def}
\end{equation}
Following Mez\H{o} and Baricz~\cite{MezoBaricz2017}, we denote it here by
\(\rW_r\).  This calligraphic notation prevents confusion with the complex
branch symbol \(\W_k\) for the ordinary Lambert function.  When \(r\ge0\),
\(h_r'(u)=\e^u(1+u)+r>0\) on \([0,\infty)\), so \(\rW_r\) is a single-valued
increasing function there.

\begin{theorem}[Exact inverse]
\label{thm:exact-inverse}
For \(z\ge0\), let \(u=\rW_1(z)\), the unique nonnegative solution of
\begin{equation}
 u\e^u+u=z.
 \label{eq:u-equation}
\end{equation}
Then
\begin{equation}
 \boxed{\quad g(z)=z-u=u\e^u.\quad}
 \label{eq:g-exact}
\end{equation}
\end{theorem}

\begin{proof}
Set \(x=g(z)\) and \(u=\W(x)\).  By \eqref{eq:Wdef}, \(x=u\e^u\).  The
defining equation \(z=f(x)=x+\W(x)\) therefore becomes
\[
 z=u\e^u+u.
\]
Monotonicity of \(h_1\) gives \(u=\rW_1(z)\).  Finally,
\(x=u\e^u=z-u\), which proves both identities.
\end{proof}

\paragraph{Lean status.}
Formal, in \path{LambertShiftInverse.lean}: with \path{Fabius.lambertShift} the map
\(f(x)=x+\W(x)\), \path{Fabius.shiftedLambertW} the root \(u\ge0\) of \(u\e^u+u=z\)
(defined as the inverse on \([0,\infty)\) of \path{Fabius.shiftedMulExp}, whose existence is
the intermediate value theorem on \([0,z]\) and whose uniqueness is strict monotonicity,
\path{Fabius.shiftedLambertW_unique}), and \path{Fabius.lambertShiftInv} the difference
\(z-u\), the theorem is the pair \path{Fabius.lambertShift_lambertShiftInv}
(\(f(z-u)=z\)) and \path{Fabius.lambertShiftInv_eq_mul_exp} (\(z-u=u\e^u\)); the
intermediate fact \(\W(g(z))=u\) is \path{Fabius.principalLambertW_lambertShiftInv}.  That
\(g\) is the two-sided inverse of \(f\) on \([0,\infty)\) is
\path{Fabius.lambertShiftInv_lambertShift} together with
\path{Fabius.lambertShift_strictMonoOn} and \path{Fabius.lambertShift_surjOn}.

\begin{corollary}[Elementary bounds]
\label{cor:bounds}
For every \(z>0\),
\begin{equation}
 z-\W(z)<g(z)<z.
 \label{eq:elementary-bounds}
\end{equation}
\end{corollary}

\begin{proof}
The right inequality follows from \(g(z)=z-u\) and \(u>0\).  Let
\(v=\W(z)\).  Since \(v\e^v=z\), one has
\(h_1(v)=v\e^v+v=z+v>z=h_1(u)\).  Monotonicity gives \(u<v\), hence
\(g(z)=z-u>z-v\).
\end{proof}

\paragraph{Lean status.}
Formal: \path{Fabius.lambertShiftInv_lt_self} and
\path{Fabius.sub_principalLambertW_lt_lambertShiftInv} in \path{LambertShiftInverse.lean},
for every \(z>0\); the second is the comparison \(u<\W(z)\)
(\path{Fabius.shiftedLambertW_lt_principalLambertW}) read off the strict monotonicity of
\(t\mapsto t\e^t\).

Differentiating \eqref{eq:u-equation} gives
\begin{equation}
 u'(z)=\frac{1}{\e^u(1+u)+1},
 \qquad
 g'(z)=\frac{\e^u(1+u)}{\e^u(1+u)+1}
      =\frac{1+u}{1+u+\e^{-u}}.
 \label{eq:gprime-exact}
\end{equation}
Thus \(0<g'(z)<1\), as expected from a strictly increasing perturbation of the
identity.

\subsection{The comparison point \texorpdfstring{\(v=\W(z)\)}{v=W(z)}}

The ordinary Lambert value \(v=\W(z)\) solves
\begin{equation}
 v\e^v=z.
 \label{eq:v-equation}
\end{equation}
Since \(u<v\), define the positive displacement
\begin{equation}
 \tau=v-u.
 \label{eq:tau-def}
\end{equation}
The inverse can then be written as
\begin{equation}
 g(z)=z-v+\tau.
 \label{eq:g-zv-tau}
\end{equation}
The entire problem is reduced to finding the small quantity \(\tau\).
Substituting \(u=v-\tau\) in \eqref{eq:u-equation} and dividing by \(\e^v\)
gives
\begin{equation}
 (v-\tau)\e^{-\tau}+(v-\tau)\e^{-v}=v.
 \label{eq:tau-before}
\end{equation}
With
\begin{equation}
 q=\e^{-v}=\frac{v}{z},
 \label{eq:q-def}
\end{equation}
this becomes the exact inverse-series equation
\begin{equation}
 q=\Phi_v(\tau),
 \qquad
 \Phi_v(\tau)=\frac{v}{v-\tau}-\e^{-\tau}.
 \label{eq:Phi-equation}
\end{equation}
Since \(\Phi_v(0)=0\) and
\begin{equation}
 \Phi_v'(0)=1+\frac1v>0,
 \label{eq:Phi-derivative}
\end{equation}
\(\Phi_v\) has a local analytic inverse at the origin.  Equation
\eqref{eq:Phi-equation} will later provide both convergence and efficient
coefficient generation.

\section{Why two asymptotic scales are necessary}
\label{sec:scales}

Let
\begin{equation}
 L=\log z,
 \qquad \ell=\log L.
 \label{eq:logs}
\end{equation}
The ordinary Lambert function has the familiar size
\(v=L-\ell+o(1)\).  Consequently
\begin{equation}
 q=\frac{v}{z}\sim\frac{L}{z}.
 \label{eq:q-size}
\end{equation}
For every fixed \(N\),
\begin{equation}
 \frac{L}{z}=o(L^{-N}).
 \label{eq:beyond-all-orders}
\end{equation}
Thus the displacement \(\tau\), which begins at order \(q\), is smaller than
every power of \(1/L\).

This yields a strict hierarchy:
\[
 z\gg L\gg\ell\gg1\gg \frac{\ell}{L}\gg\frac{\ell^2}{L^2}
 \gg\cdots\gg\frac{L}{z}\gg\frac{L^2}{z^2}\gg\cdots.
\]
The first line of asymptotics, organized only by powers of \(1/L\), determines
\(z-v\).  It cannot detect \(\tau\).  The second line, organized by powers of
\(q=v/z\), resolves the difference between \(z-v\) and the true inverse.

\begin{warningbox}
A statement such as
\[
 g(z)=z-L+\ell-\frac{\ell}{L}+O\!\left(\frac{\ell^2}{L^2}\right)
\]
is correct but incapable of revealing the correction \(\tau\asymp L/z\),
because that correction is already swallowed by every finite inverse-logarithm
remainder.  A complete answer must therefore either retain \(v=\W(z)\) as an
exact slow variable or use a nested transseries with an outer \(z^{-1}\) scale
and inner logarithmic coefficient expansions.
\end{warningbox}

\section{A differential algebra for logarithmic transseries}
\label{sec:formal-algebra}

The special function calculation is an instance of a general formal mechanism.
We develop that mechanism before returning to \(\W\).

\subsection{Slow coefficient algebras and the outer variable}

Let \(X\to+\infty\), put \(t=X^{-1}\), and introduce iterated logarithms
\begin{equation}
 L_1=\log X,
 \qquad L_2=\log L_1,
 \qquad L_3=\log L_2,\ \ldots .
 \label{eq:iterated-logs}
\end{equation}
Let \(\Ccal\) be a commutative \(\Q\)-algebra of slow coefficients containing
whatever rational functions, logarithmic polynomials, or completed inner
series are required.  It is equipped with the derivation
\begin{equation}
 \delta
 =\frac{\partial}{\partial L_1}
 +\frac{1}{L_1}\frac{\partial}{\partial L_2}
 +\frac{1}{L_1L_2}\frac{\partial}{\partial L_3}+\cdots,
 \label{eq:slow-derivation}
\end{equation}
truncated after the last iterated logarithm actually present.  Then
\(\dd c/\dd X=t\,\delta c\) for \(c\in\Ccal\).

Consider the outer formal series ring
\begin{equation}
 \Ccal[[t]]=\left\{\sum_{n\ge0}c_n t^n:c_n\in\Ccal\right\}.
 \label{eq:outer-ring}
\end{equation}
The derivative with respect to \(X\) acts as
\begin{equation}
 D:=\frac{\dd}{\dd X}
 =t\delta-t^2\frac{\partial}{\partial t}.
 \label{eq:D-operator}
\end{equation}
In particular,
\begin{equation}
 D(c t^r)=t^{r+1}(\delta-r)c.
 \label{eq:D-monomial}
\end{equation}
Every differentiation raises the outer \(t\)-valuation by one.  This simple
fact is responsible for all local finiteness below.

\begin{lemma}[Repeated differentiation]
\label{lem:repeated-D}
For \(c\in\Ccal\), \(r\ge0\), and \(k\ge0\),
\begin{equation}
 D^k(c t^r)
 =t^{r+k}\prod_{j=0}^{k-1}(\delta-r-j)c,
 \label{eq:repeated-D}
\end{equation}
where the empty product for \(k=0\) is the identity.
\end{lemma}

\begin{proof}
The case \(k=0\) is immediate, and \eqref{eq:D-monomial} is the case \(k=1\).
If the formula holds for \(k\), apply \eqref{eq:D-monomial} to the coefficient
of \(t^{r+k}\).  This appends the commuting factor \(\delta-r-k\).
\end{proof}

\subsection{Composition near the identity}

For \(A,B\in\Ccal[[t]]\), composition of \(A(X)\) with \(X+B(X)\) is defined
by the formal Taylor series
\begin{equation}
 A\bigl(X+B(X)\bigr)
 :=\sum_{k\ge0}\frac{B(X)^k}{k!}D^k A(X).
 \label{eq:formal-composition}
\end{equation}
Although \(B\) may have outer order zero, \(D^kA\) has order at least \(k\).
Hence only finitely many \(k\) contribute to each coefficient of \(t^N\), so
\eqref{eq:formal-composition} is well-defined.  This is the algebraic version
of expanding
\[
 \log(X+B)=\log X+\log(1+tB)
\]
and then expanding all subsequent iterated logarithms.

\subsection{Lagrange--B\"urmann reversion at infinity}

\begin{theorem}[Near-identity reversion]
\label{thm:general-reversion}
Let
\begin{equation}
 F(X)=X+A(X),\qquad A\in\Ccal[[t]].
 \label{eq:F-general}
\end{equation}
Then \(F\) has a unique compositional inverse of the form
\(G(X)=X+B(X)\), \(B\in\Ccal[[t]]\), and
\begin{equation}
 \boxed{
 G(X)=X+\sum_{n\ge1}\frac{(-1)^n}{n!}
 D^{n-1}\bigl(A(X)^n\bigr).}
 \label{eq:general-LB}
\end{equation}
The series is locally finite in the outer \(t\)-adic topology: its \(n\)-th
summand has valuation at least \(n-1\).
\end{theorem}

\begin{proof}
Introduce a formal parameter \(\varepsilon\) and solve
\[
 G_{\varepsilon}(X)+\varepsilon A(G_{\varepsilon}(X))=X.
\]
Writing \(G_{\varepsilon}=X+Y\), this is
\(Y=-\varepsilon A(X+Y)\).  The classical Lagrange--B\"urmann formula,
applied as a formal power series in \(\varepsilon\), gives
\[
 Y=\sum_{n\ge1}\frac{(-\varepsilon)^n}{n!}D^{n-1}A(X)^n.
\]
By \eqref{eq:D-operator}, the coefficient of \(\varepsilon^n\) has outer
valuation at least \(n-1\).  Therefore setting \(\varepsilon=1\) is legitimate:
for any prescribed outer order, only finitely many values of \(n\) occur.  The
resulting series satisfies both composition identities by formal substitution.

For uniqueness, suppose two inverses differ first at order \(t^N\).  Since
\(F'(X)=1+O(t)\), formal Taylor expansion shows that their images under \(F\)
differ by the same nonzero leading \(t^N\) block, a contradiction.
\end{proof}

The first universal terms of \eqref{eq:general-LB} are
\begin{align}
 G={}&X-A+AA'
 -A(A')^2-\frac{A^2}{2}A''\notag\\
 &+A(A')^3+\frac32A^2A'A''+\frac{A^3}{6}A'''
 +\cdots,
 \label{eq:universal-derivative-terms}
\end{align}
where primes mean \(D\).  This display is useful for hand calculations, but the
compact formula \eqref{eq:general-LB} is safer at high order.

\begin{remark}[Bell-polynomial form]
Let \(B_{m,k}\) be the partial exponential Bell polynomial.  For \(n\ge2\),
\begin{equation}
 D^{n-1}(A^n)
 =\sum_{k=1}^{n-1}(n)_{\underline{k}}A^{n-k}
 B_{n-1,k}(A',A'',\ldots,A^{(n-k)}),
 \label{eq:Bell-form}
\end{equation}
which is a direct Fa\`a di Bruno expansion of \(D^{n-1}(y^n)|_{y=A}\).
Thus every universal reversion coefficient can be generated by standard Bell
polynomial machinery.
\end{remark}

\subsection{An explicit coefficient-block formula}

Write
\begin{equation}
 A=\sum_{r\ge0}a_r t^r,
 \qquad
 A^n=\sum_{r\ge0}C_{n,r}t^r,
 \qquad
 C_{n,r}=\sum_{r_1+\cdots+r_n=r}a_{r_1}\cdots a_{r_n}.
 \label{eq:Cnr-def}
\end{equation}
If \(G=X+\sum_{N\ge0}b_Nt^N\), then \cref{lem:repeated-D,thm:general-reversion}
give the finite formula
\begin{equation}
 \boxed{
 b_N=\sum_{n=1}^{N+1}\frac{(-1)^n}{n!}
 \left[\prod_{j=0}^{n-2}
 \bigl(\delta-(N-n+1)-j\bigr)\right]
 C_{n,N-n+1}.}
 \label{eq:block-formula}
\end{equation}
The empty product at \(n=1\) is one.  The first blocks are
\begin{align}
 b_0={}&-a_0,\label{eq:b0}\\
 b_1={}&-a_1+a_0\delta a_0,\label{eq:b1}\\
 b_2={}&-a_2+(\delta-1)(a_0a_1)
 -\frac16\delta(\delta-1)(a_0^3),\label{eq:b2}\\
 b_3={}&-a_3+\frac12(\delta-2)(2a_0a_2+a_1^2)\notag\\
 &-\frac16(\delta-1)(\delta-2)(3a_0^2a_1)
 +\frac1{24}\delta(\delta-1)(\delta-2)(a_0^4).
 \label{eq:b3}
\end{align}

\begin{methodbox}
To reverse any logarithmic transseries \(F=X+A\):
choose a slow differential coefficient algebra \((\Ccal,\delta)\), expand
\(A=\sum a_rt^r\), compute the finite convolutions \(C_{n,r}\), and apply
\eqref{eq:block-formula}.  No guessing of logarithmic ansatzes is required, and
each output block depends on only finitely many input blocks.
\end{methodbox}

\section{Application of the general formula to \texorpdfstring{\(A(X)=\W(X)\)}{A(X)=W(X)}}
\label{sec:application}

The exact slow variable
\begin{equation}
 v=\W(X)
 \label{eq:v-slow}
\end{equation}
forms a particularly economical differential coefficient field.  Since
\(X=v\e^v\),
\begin{equation}
 \frac{\dd v}{\dd X}
 =\frac{v}{X(1+v)}
 =t\,\frac{v}{1+v}.
 \label{eq:dv-dX}
\end{equation}
Thus, on rational functions of \(v\), the slow derivation is
\begin{equation}
 \vartheta:=\frac{v}{1+v}\frac{\dd}{\dd v},
 \qquad D=t\vartheta-t^2\frac{\partial}{\partial t}.
 \label{eq:theta-def}
\end{equation}
In this coefficient field \(A=v\) has only one outer block, namely \(a_0=v\)
and \(a_r=0\) for \(r\ge1\).  Therefore, in the block formula
\eqref{eq:block-formula}, exactly one Lagrange term contributes to each outer
order.

\begin{theorem}[All-orders \(\W\)-anchored expansion]
\label{thm:anchored}
Let \(v=\W(z)\).  Then, formally in \(z^{-1}\),
\begin{equation}
 g(z)=z-v+\sum_{m\ge1}\frac{A_m(v)}{z^m},
 \label{eq:anchored-all}
\end{equation}
where
\begin{equation}
 \boxed{
 A_m(v)=\frac{(-1)^{m+1}}{(m+1)!}
 \prod_{j=0}^{m-1}(\vartheta-j)v^{m+1},
 \qquad
 \vartheta=\frac{v}{1+v}\frac{\dd}{\dd v}.}
 \label{eq:Am-operator}
\end{equation}
Equivalently,
\begin{equation}
 g(z)=z+\sum_{n\ge1}\frac{(-1)^n}{n!}
 \left(\frac{\dd}{\dd z}\right)^{n-1}v^n.
 \label{eq:anchored-derivative}
\end{equation}
\end{theorem}

\begin{proof}
Apply \eqref{eq:general-LB} with \(A=v\).  For a rational function \(R(v)\),
\begin{equation}
 \frac{\dd}{\dd z}\left(z^{-j}R(v)\right)
 =z^{-j-1}(\vartheta-j)R(v).
 \label{eq:shifted-theta}
\end{equation}
Repeated application to \(R(v)=v^{m+1}\) proves
\eqref{eq:Am-operator} after putting \(n=m+1\).  The \(n=1\) term is \(-v\).
\end{proof}

The first six coefficient functions are as follows:
\begin{align}
 A_1(v)={}&\frac{v^2}{1+v},
 \label{eq:A1}\\
 A_2(v)={}&\frac{v^3(v^2-2)}{2(1+v)^3},
 \label{eq:A2}\\
 A_3(v)={}&\frac{v^4(2v-1)(v^3-6v-6)}{6(1+v)^5},
 \label{eq:A3}\\
 A_4(v)={}&\frac{v^5}{24(1+v)^7}
 \bigl(6v^6-8v^5-69v^4-40v^3
       +96v^2+72v-24\bigr),
 \label{eq:A4}\\
 A_5(v)={}&\frac{v^6}{120(1+v)^9}
 \bigl(24v^8-58v^7-428v^6-151v^5
       +1320v^4\notag\\
 &\hspace{39mm}{}+1380v^3-480v^2-720v+120\bigr),
 \label{eq:A5}\\
 A_6(v)={}&\frac{v^7}{720(1+v)^{11}}
 \bigl(120v^{10}-444v^9-2920v^8+424v^7
       +16577v^6\notag\\
 &\hspace{27mm}{}+17892v^5-13530v^4-26400v^3
       -1800v^2+7200v-720\bigr).
 \label{eq:A6}
\end{align}
Substitution in \eqref{eq:anchored-all} gives a practical high-accuracy
expansion without expanding \(v\) into logarithms.

\subsection{Derivative polynomials and a triangular generator}

For symbolic computation it is useful to avoid repeated rational
differentiation from scratch.  Define polynomials \(R_{n,k}(v)\) by
\begin{equation}
 \left(\frac{\dd}{\dd z}\right)^k v^n
 =\frac{1}{z^k}\frac{v^n R_{n,k}(v)}{(1+v)^{2k-1}},
 \qquad k\ge1.
 \label{eq:Rnk-def}
\end{equation}
Then
\begin{equation}
 R_{n,1}(v)=n,
 \label{eq:Rnk-initial}
\end{equation}
and differentiation of \eqref{eq:Rnk-def} gives
\begin{align}
 R_{n,k+1}(v)={}&v(1+v)R'_{n,k}(v)\notag\\
 &+\bigl(n(1+v)-(2k-1)v-k(1+v)^2\bigr)R_{n,k}(v).
 \label{eq:Rnk-recurrence}
\end{align}
Consequently,
\begin{equation}
 A_m(v)=\frac{(-1)^{m+1}}{(m+1)!}
 \frac{v^{m+1}R_{m+1,m}(v)}{(1+v)^{2m-1}}.
 \label{eq:Am-Rnk}
\end{equation}
This two-index recurrence generates the diagonal needed for the inverse and,
at the same time, all derivatives of all powers of \(\W\).

\begin{proposition}[Integer-polynomial structure]
\label{prop:Pm-structure}
For every \(m\ge1\), there is a polynomial \(P_m\in\mathbb{Z}[v]\) such that
\begin{equation}
 A_m(v)=\frac{v^{m+1}P_m(v)}{m!(1+v)^{2m-1}}.
 \label{eq:Am-Pm}
\end{equation}
Moreover,
\begin{equation}
 \deg P_m=2m-2,
 \qquad [v^{2m-2}]P_m=(m-1)!,
 \qquad P_m(0)=(-1)^{m+1}m!.
 \label{eq:Pm-data}
\end{equation}
In particular,
\begin{equation}
 A_m(v)\sim\frac{v^m}{m}
 \qquad (v\to+\infty).
 \label{eq:Am-leading}
\end{equation}
\end{proposition}

\begin{proof}
Fix a positive integer \(n\).  Since \(R_{n,1}=n\) and every operation on the
right-hand side of \eqref{eq:Rnk-recurrence} preserves integer coefficients
and divisibility by \(n\), one has
\(R_{n,k}\in n\mathbb Z[v]\) for every \(k\ge1\).  Consequently
\[
 P_m(v)=\frac{(-1)^{m+1}}{m+1}R_{m+1,m}(v)
\]
is an integer polynomial, and \eqref{eq:Am-Pm} follows from
\eqref{eq:Am-Rnk}.

It remains to determine its extremal coefficients.  If
\(d_k=\deg R_{n,k}\) and \(r_k\) is the leading coefficient, then
\eqref{eq:Rnk-recurrence} shows, inductively, that
\[
 d_1=0,\qquad d_{k+1}=d_k+2,
 \qquad r_{k+1}=-k r_k.
\]
Indeed, the unique term of degree \(d_k+2\) is
\(-kv^2R_{n,k}\).  Hence
\[
 \deg R_{n,k}=2k-2,
 \qquad [v^{2k-2}]R_{n,k}=(-1)^{k-1}(k-1)!n.
\]
On the diagonal \(n=m+1,k=m\), multiplication by
\((-1)^{m+1}/(m+1)\) gives
\(\deg P_m=2m-2\) and leading coefficient \((m-1)!\).

Finally, evaluating \eqref{eq:Rnk-recurrence} at \(v=0\) gives
\[
 R_{n,k+1}(0)=(n-k)R_{n,k}(0),
 \qquad R_{n,1}(0)=n.
\]
Therefore \(R_{n,k}(0)=n!/(n-k)!\) for \(1\le k\le n\), and
\[
 P_m(0)=\frac{(-1)^{m+1}}{m+1}(m+1)!
       =(-1)^{m+1}m!.
\]
Substitution in \eqref{eq:Am-Pm} now gives
\(A_m(v)\sim v^m/m\), completing the proof.
\end{proof}

\section{Direct reversion in the exponentially small parameter}
\label{sec:q-reversion}

The formal Lagrange expansion already gives the answer.  The exact equation
\eqref{eq:Phi-equation} gives more: a convergent series for all sufficiently
large \(z\), a second coefficient recurrence, and a natural remainder scale.

\subsection{Lagrange inversion of \texorpdfstring{\(\Phi_v\)}{Phi-v}}

Let
\begin{equation}
 \tau(q;v)=\sum_{n\ge1}c_n(v)q^n
 \label{eq:tau-c-series}
\end{equation}
be the local inverse of \(q=\Phi_v(\tau)\).  Lagrange inversion gives
\begin{equation}
 \boxed{
 c_n(v)=\frac1n[\tau^{n-1}]
 \left(\frac{\tau}{\Phi_v(\tau)}\right)^n
 =\frac1{n!}\left.
 \frac{\dd^{n-1}}{\dd\tau^{n-1}}
 \left(\frac{\tau}{\Phi_v(\tau)}\right)^n
 \right|_{\tau=0}.}
 \label{eq:cn-Lagrange}
\end{equation}
Since \(q=v/z\), comparison with \eqref{eq:anchored-all} yields
\begin{equation}
 A_n(v)=v^n c_n(v).
 \label{eq:An-cn}
\end{equation}
For example,
\begin{align}
 c_1(v)={}&\frac{v}{1+v},\label{eq:c1}\\
 c_2(v)={}&\frac{v(v^2-2)}{2(1+v)^3},\label{eq:c2}\\
 c_3(v)={}&\frac{v(2v-1)(v^3-6v-6)}{6(1+v)^5}.
 \label{eq:c3}
\end{align}

Expanding \(\Phi_v\) itself,
\begin{equation}
 \Phi_v(\tau)=\sum_{k\ge1}\phi_k(v)\tau^k,
 \qquad
 \phi_k(v)=\frac1{v^k}-\frac{(-1)^k}{k!},
 \label{eq:phi-k}
\end{equation}
produces a purely triangular recurrence.  If
\(C(q)=\sum_{n\ge1}c_nq^n\), then the identity
\(\sum_{k\ge1}\phi_k C(q)^k=q\) gives
\begin{align}
 c_1={}&\frac1{\phi_1},
 \label{eq:cn-rec-first}\\
 c_n={}&-\frac1{\phi_1}
 \sum_{k=2}^{n}\phi_k
 \sum_{\substack{j_1+\cdots+j_k=n\\j_i\ge1}}
 c_{j_1}\cdots c_{j_k},
 \qquad n\ge2.
 \label{eq:cn-rec}
\end{align}
All arithmetic is rational in \(v\).

\subsection{Convergence and a natural remainder theorem}

\begin{theorem}[Convergent outer expansion for large \(z\)]
\label{thm:q-convergence}
There exist constants \(v_0>0\), \(\rho>0\), and \(M_N>0\) such that for all
real \(v\ge v_0\), the inverse \(\tau(q;v)=\Phi_v^{-1}(q)\) is analytic for
\(|q|<\rho\), with \(\rho\) independent of \(v\).  Consequently, for
\(z=v\e^v\) and every fixed \(N\),
\begin{equation}
 g(z)=z-v+\sum_{m=1}^{N}\frac{A_m(v)}{z^m}
 +O_N\!\left(\left(\frac{v}{z}\right)^{N+1}\right),
 \qquad z\to+\infty.
 \label{eq:q-remainder}
\end{equation}
For all sufficiently large \(z\), the infinite series
\eqref{eq:anchored-all} converges to \(g(z)\).
\end{theorem}

\begin{proof}
Put \(s=1/v\) and write
\begin{equation}
 \Phi_s(\tau)=\frac1{1-s\tau}-\e^{-\tau}.
 \label{eq:Phi-s}
\end{equation}
This is analytic in \((s,\tau)\) near \((0,0)\), and
\[
 \Phi_s(0)=0,
 \qquad
 \partial_{\tau}\Phi_s(0)=1+s.
\]
The analytic inverse function theorem with parameter therefore gives a
bidisc \(|s|<s_0\), \(|q|<\rho\) on which the inverse \(\tau=T(q,s)\) is
analytic.  Compactness after slightly shrinking the bidisc makes \(\rho\)
uniform for \(0\le s\le s_0/2\), equivalently \(v\ge v_0\).

Now the actual parameter is \(q=\e^{-v}\), which is eventually smaller than
\(\rho\).  Taylor's theorem in \(q\), uniformly for the indicated values of
\(s\), gives a remainder bounded by \(M_N|q|^{N+1}\).  Substituting
\(q=v/z\), \(\tau=g-(z-v)\), and \(A_m=v^mc_m\) proves
\eqref{eq:q-remainder}.  The same analytic inverse theorem proves convergence
of the full Taylor series.
\end{proof}

\begin{remark}
The proof is intentionally local.  It neither requires nor asserts a simple
closed form for the exact radius of convergence as a function of \(v\).  For
asymptotics, a uniform positive radius is sufficient because the actual
argument \(q=\e^{-v}\) tends to zero exponentially.
\end{remark}

\subsection{A bivariate re-expansion and harmonic numbers}

The parameter form \eqref{eq:Phi-s} permits another useful organization:
expand first in \(s=1/v\), while retaining the entire dependence on
\(q\).  Write
\begin{equation}
 T(q,s)=\sum_{j\ge0}s^jT_j(q),
 \qquad E=1-q.
 \label{eq:Tjs}
\end{equation}
At \(s=0\), \(q=1-\e^{-T_0}\), so
\begin{equation}
 T_0(q)=-\log(1-q).
 \label{eq:T0}
\end{equation}
Coefficient comparison in
\(1/(1-sT)-\e^{-T}=q\) gives successively
\begin{align}
 T_1(q)={}&\frac{\log(1-q)}{1-q},
 \label{eq:T1}\\
 T_2(q)={}&
 \frac{-\log(1-q)+\tfrac12\log^2(1-q)}{(1-q)^2}
 -\frac{\log^2(1-q)}{1-q}.
 \label{eq:T2}
\end{align}
Thus
\begin{align}
 \tau={}&-\log(1-q)
 +\frac1v\frac{\log(1-q)}{1-q}\notag\\
 &+\frac1{v^2}\left(
 \frac{-\log(1-q)+\tfrac12\log^2(1-q)}{(1-q)^2}
 -\frac{\log^2(1-q)}{1-q}\right)
 +O(v^{-3}).
 \label{eq:tau-resummed}
\end{align}
This formula resums infinitely many powers of \(q\) at each displayed order in
\(1/v\).

The construction is triangular to every order.  Indeed, at order \(s^j\), the
unknown \(T_j\) occurs only through the linear term \((1-q)T_j\) generated by
\(-\e^{-T}\).  Hence \(T_j\) is obtained by dividing an explicit expression in
\(T_0,\ldots,T_{j-1}\) by \(1-q\).  Inductively,
\begin{equation}
 T_j(q)\in(1-q)^{-j}\Q[1-q,\log(1-q)].
 \label{eq:Tj-class}
\end{equation}

Let
\begin{equation}
 H_n=\sum_{k=1}^{n}\frac1k,
 \qquad
 H_n^{(2)}=\sum_{k=1}^{n}\frac1{k^2}.
 \label{eq:harmonic}
\end{equation}
The generating identities
\begin{align}
 -\log(1-q)={}&\sum_{n\ge1}\frac{q^n}{n},\notag\\
 \frac{\log(1-q)}{1-q}={}&-\sum_{n\ge1}H_nq^n,\notag\\
 \frac{\log^2(1-q)}{1-q}={}&
 \sum_{n\ge1}\bigl(H_n^2-H_n^{(2)}\bigr)q^n
 \label{eq:harmonic-generating}
\end{align}
then imply the uniform large-\(v\) coefficient expansion
\begin{equation}
 \boxed{
 c_n(v)=\frac1n-\frac{H_n}{v}
 +\frac{H_n+\frac{n-1}{2}
   (H_n^2-H_n^{(2)})}{v^2}
 +O_n(v^{-3}).}
 \label{eq:cn-harmonic}
\end{equation}
The leading limit \(c_n(v)\to1/n\) reflects the limiting inverse
\(T_0(q)=-\log(1-q)\).

\section{The complete inverse-logarithm expansion}
\label{sec:pure-log}

The ordinary Lambert function has the following expansion, which converges
for all sufficiently large positive \(z\)~\cite{CorlessEtAl1996,DLMF413}:
\begin{equation}
 v=\W(z)=L-\ell+\sum_{n\ge1}\frac{P_n(\ell)}{L^n},
 \label{eq:W-log-expansion}
\end{equation}
with \(P_n\) given by \eqref{eq:Pn-stirling}.  The first polynomials are
\begin{align}
 P_1(y)={}&y,\notag\\
 P_2(y)={}&\frac{y(y-2)}2,\notag\\
 P_3(y)={}&\frac{y(2y^2-9y+6)}6,\notag\\
 P_4(y)={}&\frac{y(3y^3-22y^2+36y-12)}{12},\notag\\
 P_5(y)={}&\frac{y(12y^4-125y^3+350y^2-300y+60)}{60},\notag\\
 P_6(y)={}&\frac{y(20y^5-274y^4+1125y^3-1700y^2+900y-120)}{120}.
 \label{eq:P1-P6}
\end{align}

\begin{theorem}[Pure logarithmic asymptotics of the inverse]
\label{thm:pure-log}
For every fixed \(N\ge0\),
\begin{equation}
 g(z)=z-L+\ell-\sum_{n=1}^{N}\frac{P_n(\ell)}{L^n}
 +O_N\!\left(\frac{\ell^{N+1}}{L^{N+1}}\right),
 \qquad z\to+\infty.
 \label{eq:pure-log-truncated}
\end{equation}
Equivalently, \(g(z)\) and \(z-\W(z)\) have the same complete asymptotic
expansion in the inverse-logarithm scale.
\end{theorem}

\begin{proof}
By \cref{thm:q-convergence},
\[
 g(z)-(z-v)=O(v/z)=O(L/z).
\]
Equation \eqref{eq:beyond-all-orders} shows that \(L/z=o(L^{-M})\) for every
fixed \(M\).  Therefore the difference is beyond all orders in \(1/L\), and
substitution of the standard Lambert expansion \eqref{eq:W-log-expansion}
proves the result.
\end{proof}

\paragraph{Lean status.}
Partial.  The Lambert-side input, the leading two terms of \eqref{eq:W-log-expansion}, is
formal in \path{PrincipalLambertWAtTop.lean}: the identity \(\W(x)+\log\W(x)=\log x\)
(\path{Fabius.principalLambertW_add_log}), the bracket \(L-\ell\le\W(x)\le L\) for \(x\ge\e\)
(\path{Fabius.principalLambertW_log_bracket}), \(\W\sim\log\) at \(+\infty\)
(\path{Fabius.principalLambertW_isEquivalent_log}) and the two-term statement
\(\W(x)-(L-\ell)\to0\) (\path{Fabius.tendsto_principalLambertW_sub_log_sub_log_log}).  The
higher terms \(P_n(\ell)/L^n\), and the beyond-all-orders step
\(g(z)-(z-\W(z))=O(L/z)\) that transfers the expansion from \(z-\W(z)\) to \(g\), are not
formalized.

In explicit form,
\begin{align}
 g(z)\sim{}&z-L+\ell-\frac{\ell}{L}
 -\frac{\ell(\ell-2)}{2L^2}
 -\frac{\ell(2\ell^2-9\ell+6)}{6L^3}\notag\\
 &-\frac{\ell(3\ell^3-22\ell^2+36\ell-12)}{12L^4}\notag\\
 &-\frac{\ell(12\ell^4-125\ell^3+350\ell^2-300\ell+60)}{60L^5}\notag\\
 &-\frac{\ell(20\ell^5-274\ell^4+1125\ell^3-1700\ell^2
   +900\ell-120)}{120L^6}+\cdots.
 \label{eq:pure-log-explicit}
\end{align}

\begin{warningbox}[title={What the complete slow expansion still does not know}]
Even if the series \eqref{eq:W-log-expansion} is summed to the exact value
\(v=\W(z)\), it gives only \(z-v\), not \(g(z)\).  The missing positive amount
is
\[
 \tau=g(z)-(z-v)
 =\frac{v^2}{z(1+v)}+O\!\left(\frac{v^2}{z^2}\right)
 \sim\frac{\log z}{z}.
\]
Thus the generalized inverse contains information that is invisible to the
entire rank-one inverse-logarithm expansion.  This is a concrete example of why
asymptotic equality to all algebraic orders does not imply equality of
functions.
\end{warningbox}

\section{The full two-level logarithmic transseries}
\label{sec:full-transseries}

The most informative answer is the nested expansion
\begin{equation}
 \boxed{
 g(z)\sim_{\Tcal}z-v+\sum_{m\ge1}\frac{A_m(v)}{z^m},
 \qquad
 v\sim L-\ell+\sum_{n\ge1}\frac{P_n(\ell)}{L^n}.}
 \label{eq:nested-transseries}
\end{equation}
The symbol \(\sim_{\Tcal}\) emphasizes a lexicographically ordered transseries:
first by the outer exponent of \(z^{-1}\), and inside each outer block by the
inverse powers of \(L\), with polynomial dependence on \(\ell\).  Keeping
\(v=\W(z)\) exact is usually computationally superior; expanding \(v\) is
useful when only elementary logarithms are desired.

\subsection{Leading inner structure of every outer block}

Combining \eqref{eq:An-cn} and \eqref{eq:cn-harmonic} gives
\begin{equation}
 A_m(v)=\frac{v^m}{m}-H_m v^{m-1}
 +d_m v^{m-2}+O_m(v^{m-3}),
 \label{eq:Am-v-large}
\end{equation}
where
\begin{equation}
 d_m=H_m+\frac{m-1}{2}\bigl(H_m^2-H_m^{(2)}\bigr).
 \label{eq:dm-def}
\end{equation}
Using
\(v=L-\ell+\ell/L+O(\ell^2/L^2)\), one obtains a uniform formula for the
first three inner layers of every outer block:
\begin{align}
 A_m(\W(z))={}&\frac{L^m}{m}-(\ell+H_m)L^{m-1}\notag\\
 &+\left[
 \frac{m-1}{2}\ell^2+\bigl(1+(m-1)H_m\bigr)\ell+d_m
 \right]L^{m-2}
 +O_m(L^{m-3}\ell^3).
 \label{eq:Am-log-inner}
\end{align}
For the first three blocks this reads
\begin{align}
 A_1(\W(z))={}&L-\ell-1+\frac{\ell+1}{L}
 +O\!\left(\frac{\ell^2}{L^2}\right),
 \label{eq:A1-log}\\
 A_2(\W(z))={}&\frac{L^2}{2}-\left(\ell+\frac32\right)L
 +\frac{\ell^2}{2}+\frac52\ell+2
 +O\!\left(\frac{\ell^2}{L}\right),
 \label{eq:A2-log}\\
 A_3(\W(z))={}&\frac{L^3}{3}-\left(\ell+\frac{11}{6}\right)L^2
 +\left(\ell^2+\frac{14}{3}\ell+\frac{23}{6}\right)L
 +O(\ell^3).
 \label{eq:A3-log}
\end{align}
Hence the first beyond-all-orders sector begins as
\begin{equation}
 \frac{A_1(v)}{z}
 =\frac{L-\ell-1}{z}
 +\frac{\ell+1}{zL}
 +O\!\left(\frac{\ell^2}{zL^2}\right),
 \label{eq:first-exponential-sector}
\end{equation}
and the next sector begins with \(L^2/(2z^2)\).

\begin{remark}[Why the leading coefficients are logarithmic]
At fixed \(m\), \(A_m(v)/z^m\sim q^m/m\).  Summing only these leading pieces
over all \(m\) gives
\[
 \sum_{m\ge1}\frac{q^m}{m}=-\log(1-q),
\]
which is exactly the \(s=0\) resummation \(T_0(q)\) in
\eqref{eq:T0}.  The harmonic correction \(-H_m/v\) in \eqref{eq:cn-harmonic}
resums to \(T_1(q)/v\).  Thus the rational coefficient blocks, the harmonic
numbers, and the elementary logarithm \(-\log(1-q)\) are three views of the
same bivariate inverse.
\end{remark}

\subsection{A transseries ordering convention}

A practical truncation is specified by two cutoffs.  Choose an outer order
\(M\) and, for each \(0\le m\le M\), an inner order \(N_m\).  Then:
\begin{enumerate}
\item expand the slow block \(-v\) through \(L^{-N_0}\);
\item for \(1\le m\le M\), expand \(A_m(v)\) through
      \(L^{m-N_m}\), then multiply by \(z^{-m}\);
\item retain the unexpanded form \(A_m(\W(z))/z^m\) whenever numerical
      evaluation of \(\W\) is available.
\end{enumerate}
A single scalar ``order'' is insufficient because an error such as
\(L^{-100}\) is still much larger than the leading \(L/z\) sector.  This is
not a defect; it is the defining feature of a rank-two asymptotic scale.

\section{Algorithms for general logarithmic-transseries reversion}
\label{sec:algorithms}

The closed formula \eqref{eq:general-LB} is conceptually definitive, but
several algorithms are useful in different computational regimes.

\subsection{Algorithm I: direct Lagrange--B\"urmann blocks}

\begin{algorithm}[Explicit block construction]
\label{alg:LB-block}
Given \(A=\sum_{r=0}^{N}a_rt^r+O(t^{N+1})\):
\begin{enumerate}
\item form \(C_{n,r}=[t^r]A^n\) for \(1\le n\le N+1\) and
      \(0\le r\le N-n+1\);
\item apply the falling differential operator in \eqref{eq:block-formula};
\item sum over \(n\) to obtain \(b_N\), and repeat for all lower blocks if
      they have not already been computed.
\end{enumerate}
The output is the inverse \(G=X+\sum_{r=0}^{N}b_rt^r+O(t^{N+1})\).
\end{algorithm}

This method gives transparent closed formulas and is ideal when exact
coefficient identities matter.  Its disadvantage is repeated formation of
powers \(A^n\) if implemented naively.

\subsection{Algorithm II: triangular residual cancellation}

Suppose an approximation \(G_{<N}\) has been constructed with
\begin{equation}
 F(G_{<N})-X=r_Nt^N+O(t^{N+1}).
 \label{eq:triangular-residual}
\end{equation}
Because \(F'(G_{<N})=1+O(t)\), the correction
\begin{equation}
 G_{<N+1}=G_{<N}-r_Nt^N
 \label{eq:triangular-correction}
\end{equation}
raises the residual order by at least one.  Starting with \(G_{<0}=X\), this
builds the inverse one block at a time.  The first step automatically gives
\(b_0=-a_0\).

\begin{algorithm}[Triangular reversion]
\label{alg:triangular}
For \(N=0,1,2,\ldots\): compose \(F\) with the current truncation, extract the
first nonzero residual block \(r_N\), subtract \(r_Nt^N\) from the current
inverse, and truncate all intermediate expressions above the target order.
\end{algorithm}

The triangular method is often the simplest robust implementation.  It also
provides a residual certificate at every stage.

\subsection{Algorithm III: fixed-point iteration}

The inverse equation is
\begin{equation}
 G=X-A(G).
 \label{eq:fixed-point-equation}
\end{equation}
Starting with \(G_0=X\), define
\begin{equation}
 G_{k+1}=X-A(G_k).
 \label{eq:fixed-point-iteration}
\end{equation}
If two approximants first differ at order \(t^N\), composition with \(A\)
raises that difference to order at least \(t^{N+1}\), because \(A'=O(t)\).
Thus each iteration gains at least one outer block.  This is a formal
contraction in the \(t\)-adic metric.

For the present problem,
\begin{equation}
 G_{k+1}(z)=z-\W(G_k(z)).
 \label{eq:W-fixed-point}
\end{equation}
Numerically this is already contractive for every positive argument, since
\(0<\W'(x)\le1\), and it becomes rapidly contractive as \(x\to\infty\).

\subsection{Algorithm IV: Newton reversion}

Newton's method applies directly to \(F(G)-X=0\):
\begin{equation}
 G_{\mathrm{new}}
 =G-\frac{F(G)-X}{F'(G)}.
 \label{eq:Newton-general}
\end{equation}
In the filtered transseries algebra, reciprocal computation is legitimate
because \(F'(G)=1+O(t)\) is a unit.  If the residual has outer order \(N\), the
new residual has at least approximately twice that precision; for a slow
perturbation, \(F''=O(t^2)\) makes the gain even stronger.

For \(F(x)=x+\W(x)\),
\begin{equation}
 G_{\mathrm{new}}
 =G-\frac{G+\W(G)-z}
 {1+\W(G)/(G(1+\W(G)))}.
 \label{eq:Newton-W}
\end{equation}
A low-order transseries seed followed by one Newton step is an effective
hybrid for high-precision numerical evaluation.

\subsection{Algorithm V: direct inverse in a problem-adapted small parameter}

Whenever an exact substitution reduces the problem to an analytic equation
\(q=\Phi(\tau;\lambda)\) with \(\Phi(0;\lambda)=0\) and
\(\partial_\tau\Phi(0;\lambda)\ne0\), one should consider reversing \(\Phi\)
directly.  In the present case, \(q=v/z\) is exponentially small and
\eqref{eq:cn-rec} is cheaper than generic composition.  This method also
exposes actual convergence, which a purely formal outer calculation does not
automatically prove.

\begin{summarybox}
The five methods are mathematically equivalent but computationally distinct.
Lagrange--B\"urmann gives closed coefficients; triangular cancellation is the
simplest certified block engine; fixed-point iteration is structurally
transparent; Newton is fastest at high precision; a problem-adapted inverse
such as \(q=\Phi_v(\tau)\) can reveal convergence and hidden resummations.
\end{summarybox}

\section{Analytic transfer and residual certificates}
\label{sec:certification}

A formal transseries identity does not, by itself, state an analytic error
bound.  For monotone real inverse problems, the residual supplies an elementary
and often sharp bridge.

\begin{theorem}[Residual-to-error transfer]
\label{thm:residual-general}
Let \(F\) be continuously differentiable and strictly increasing on an interval
containing the exact inverse value \(G(X)\) and an approximation \(H(X)\).  If
\begin{equation}
 0<m\le F'(y)\le M
 \label{eq:derivative-bounds}
\end{equation}
throughout the segment between \(G(X)\) and \(H(X)\), and
\(R=F(H)-X\), then
\begin{equation}
 \frac{|R|}{M}\le |H-G|\le\frac{|R|}{m}.
 \label{eq:residual-error-bound}
\end{equation}
\end{theorem}

\begin{proof}
The mean value theorem gives
\(R=F(H)-F(G)=F'(\xi)(H-G)\) for some intermediate \(\xi\).
Apply \eqref{eq:derivative-bounds}.
\end{proof}

\paragraph{Lean status.}
The instance used below, \(F=f\), \(m=1\), \(M=2\), is formal in
\path{LambertShiftInverse.lean}: the derivative bracket \eqref{eq:Fprime-bound} is
\path{Fabius.one_le_deriv_lambertShift} and \path{Fabius.deriv_lambertShift_le_two} (from
\path{Fabius.lambertShift_hasDerivAt} and the bound \(\W'\le1\),
\path{Fabius.inv_exp_mul_add_one_le_one}), and the mean-value transfer is
\path{Fabius.abs_sub_lambertShiftInv_le_abs_residual} (\(|H-G|\le|R|\)) and
\path{Fabius.abs_residual_le_two_mul_abs_sub_lambertShiftInv} (\(|R|\le2|H-G|\)), for every
\(z\ge0\) and every approximation \(H\ge0\).  The general \((m,M)\) statement is not
separately formalized.

For \(F(x)=x+\W(x)\), \eqref{eq:Wprime} gives
\begin{equation}
 0<\W'(x)=\frac{\e^{-\W(x)}}{1+\W(x)}\le1,
 \qquad x\ge0.
 \label{eq:Wprime-bound}
\end{equation}
Therefore
\begin{equation}
 1\le F'(x)\le2.
 \label{eq:Fprime-bound}
\end{equation}
For every approximation \(H\ge0\), we obtain the particularly simple
certificate
\begin{equation}
 \boxed{
 \frac12|H+\W(H)-z|
 \le |H-g(z)|
 \le |H+\W(H)-z|.}
 \label{eq:special-residual-certificate}
\end{equation}
Thus every symbolic truncation can be checked numerically without first
computing the exact inverse.

For an outer truncation
\begin{equation}
 H_N(z)=z-v+\sum_{m=1}^{N}\frac{A_m(v)}{z^m},
 \label{eq:HN-def}
\end{equation}
\cref{thm:q-convergence} gives
\begin{equation}
 H_N-g=O_N(q^{N+1}).
 \label{eq:HN-error}
\end{equation}
The residual certificate independently verifies the same order and catches
coefficient or truncation mistakes in an implementation.

\section{Numerical behavior}
\label{sec:numerics}

The following table compares the exact inverse with successive
\(\W\)-anchored truncations.  The exact reference value was obtained by solving
\(u\e^u+u=z\) at high precision and using \(g=z-u\).  The column \(M=0\) is
\(z-\W(z)\); for \(M\ge1\), the approximation is \(H_M\) from
\eqref{eq:HN-def}.  All entries after the \(q\) column are absolute errors.

\begin{table}[htbp]
\centering
\resizebox{\textwidth}{!}{%
\begin{tabular}{@{}rrrrrrr@{}}
\toprule
\(z\) & \(q=\W(z)/z\) & \(M=0\) & \(M=1\) & \(M=2\) & \(M=4\) & \(M=6\)\\
\midrule
\(10\) & \(1.7455280\!\times10^{-1}\) &
\(1.1202183\!\times10^{-1}\) &
\(1.0461658\!\times10^{-3}\) &
\(2.9896578\!\times10^{-4}\) &
\(6.8888567\!\times10^{-7}\) &
\(1.1896934\!\times10^{-8}\)\\
\(10^2\) & \(3.3856301\!\times10^{-2}\) &
\(2.6355095\!\times10^{-2}\) &
\(2.1861948\!\times10^{-4}\) &
\(9.4974943\!\times10^{-7}\) &
\(7.9460859\!\times10^{-10}\) &
\(5.5483129\!\times10^{-14}\)\\
\(10^4\) & \(7.2318460\!\times10^{-4}\) &
\(6.3550313\!\times10^{-4}\) &
\(1.7057929\!\times10^{-7}\) &
\(5.3410695\!\times10^{-11}\) &
\(4.0832057\!\times10^{-18}\) &
\(2.9566636\!\times10^{-25}\)\\
\(10^8\) & \(1.5668997\!\times10^{-7}\) &
\(1.4728989\!\times10^{-7}\) &
\(1.0113357\!\times10^{-14}\) &
\(8.8743507\!\times10^{-22}\) &
\(8.3618466\!\times10^{-36}\) &
\(8.2973506\!\times10^{-50}\)\\
\bottomrule
\end{tabular}%
}
\caption{Absolute error of the \(\W\)-anchored expansion.}
\label{tab:numerical-errors}
\end{table}

Several features are visible.
\begin{itemize}
\item The error of \(z-v\) is not a defect in the ordinary Lambert expansion;
      it is the omitted transseries sector \(\tau\sim q\).
\item Each additional outer block usually reduces the error by approximately
      a factor of \(q\), in agreement with \eqref{eq:q-remainder}.
\item Even at \(z=10\), where \(q\approx0.175\) is not extremely small, six
      correction blocks give about eight decimal digits in absolute error.
\item At large \(z\), evaluating \(v=\W(z)\) once and applying rational
      corrections is dramatically more efficient than expanding all slow
      logarithmic subblocks to comparable accuracy.
\end{itemize}

\section{Parameter generalizations}
\label{sec:generalizations}

\subsection{The family \texorpdfstring{\(x+r\W(x)\)}{x+rW(x)}}

Let
\begin{equation}
 f_r(x)=x+r\W(x),\qquad r\ge0,
 \label{eq:fr-def}
\end{equation}
and let \(g_r=f_r^{-1}\).  Setting \(u=\W(g_r(z))\) gives
\begin{equation}
 u\e^u+ru=z,
 \label{eq:r-u-equation}
\end{equation}
so
\begin{equation}
 g_r(z)=z-r\rW_r(z).
 \label{eq:gr-exact}
\end{equation}
The near-identity formula yields
\begin{equation}
 g_r(z)=z+\sum_{n\ge1}\frac{(-r)^n}{n!}
 \left(\frac{\dd}{\dd z}\right)^{n-1}\W(z)^n.
 \label{eq:gr-LB}
\end{equation}
In terms of the coefficient functions already computed,
\begin{equation}
 \boxed{
 g_r(z)=z-rv+\sum_{m\ge1}\frac{r^{m+1}A_m(v)}{z^m},
 \qquad v=\W(z).}
 \label{eq:gr-anchored}
\end{equation}
Equivalently, for \(r\ne0\),
\begin{equation}
 \rW_r(z)=v-\sum_{m\ge1}\frac{r^mA_m(v)}{z^m}.
 \label{eq:rW-expansion}
\end{equation}
The direct equation becomes \(rq=\Phi_v(\tau)\), so the same coefficient
functions \(c_n(v)\) apply after replacing \(q\) by \(rq\).

\subsection{The scaled family \texorpdfstring{\(x+\alpha\W(\beta x)\)}{x+alpha W(beta x)}}

Let \(\alpha,\beta>0\) and
\begin{equation}
 F_{\alpha,\beta}(x)=x+\alpha\W(\beta x).
 \label{eq:scaled-family}
\end{equation}
If \(u=\W(\beta x)\), then \(x=u\e^u/\beta\), and the inverse equation is
\begin{equation}
 u\e^u+\alpha\beta u=\beta z.
 \label{eq:scaled-u}
\end{equation}
Hence
\begin{equation}
 F_{\alpha,\beta}^{-1}(z)
 =z-\alpha\rW_{\alpha\beta}(\beta z).
 \label{eq:scaled-exact}
\end{equation}
Putting \(v=\W(\beta z)\), the same universal rational functions give
\begin{equation}
 \boxed{
 F_{\alpha,\beta}^{-1}(z)
 =z-\alpha v+\sum_{m\ge1}
 \frac{\alpha^{m+1}A_m(v)}{z^m}.}
 \label{eq:scaled-expansion}
\end{equation}
The parameter \(\beta\) enters only through the anchor \(v=\W(\beta z)\).
The cancellation follows from
\[
 \frac{\dd}{\dd z}\W(\beta z)=\frac{v}{z(1+v)}.
\]

\subsection{General slowly varying perturbations}

The same machinery applies to
\begin{equation}
 F(X)=X+A(X),
 \qquad A(X)=o(X),
 \label{eq:slow-general}
\end{equation}
whenever derivatives of \(A\) lower the asymptotic order by one outer power.
Examples include finite combinations of iterated logarithms, rational
functions of \(\log X\), completed power-logarithmic series, and functions
living in a differential Hardy field or a suitable logarithmic-exponential
transseries field.  The essential structural hypotheses are:
\begin{enumerate}
\item a complete outer filtration in which \(A\) has order zero relative to
      \(t=X^{-1}\);
\item a derivation \(D\) that raises outer valuation;
\item locally finite Taylor composition by \(X+B\);
\item invertibility of the leading derivative, here the unit
      \(F'=1+O(t)\).
\end{enumerate}
Under these conditions, \eqref{eq:general-LB}, \eqref{eq:block-formula}, and the
residual algorithms remain valid verbatim.

\section{Further deductions and structural observations}
\label{sec:further}

\subsection{An expansion of the generalized Lambert value itself}

Since \(u=\rW_1(z)=z-g(z)\), \cref{thm:anchored} immediately gives
\begin{equation}
 \rW_1(z)=v-\sum_{m\ge1}\frac{A_m(v)}{z^m}.
 \label{eq:rW1-expansion}
\end{equation}
Thus
\begin{align}
 \rW_1(z)={}&v-\frac{v^2}{z(1+v)}
 -\frac{v^3(v^2-2)}{2z^2(1+v)^3}\notag\\
 &-\frac{v^4(2v-1)(v^3-6v-6)}{6z^3(1+v)^5}
 +O\!\left(\left(\frac{v}{z}\right)^4\right).
 \label{eq:rW1-first}
\end{align}
This refines the leading generalized-Lambert asymptotics while keeping the
ordinary Lambert function as an exact comparison solution.

\subsection{The correction as a perturbation of Wright omega}

The equations for \(u\) and \(v\) may also be written
\begin{align}
 v+\log v={}&L,\label{eq:v-omega}\\
 u+\log u+\log(1+\e^{-u})={}&L.
 \label{eq:u-omega}
\end{align}
The first equation says \(v=\omega(L)\), where \(\omega\) is the Wright omega
function.  The last logarithm in \eqref{eq:u-omega} is exponentially small in
\(u\); it is precisely the source of the missing transseries sector.  This
provides a second conceptual explanation for why the complete inverse-logarithm
expansion of \(u\) agrees with that of \(v\), while the exact functions differ.

\subsection{Closure under differentiation}

The anchored coefficient field \(\Q(v)\) is stable under the slow derivation
\(\vartheta\).  Therefore the complete transseries
\eqref{eq:anchored-all} may be differentiated term by term in its domain of
convergence.  For example,
\begin{equation}
 g'(z)=1-\frac{v}{z(1+v)}
 +\frac1{z^2}(\vartheta-1)A_1(v)
 +\frac1{z^3}(\vartheta-2)A_2(v)+\cdots.
 \label{eq:gprime-series}
\end{equation}
This must agree with the exact formula \eqref{eq:gprime-exact} after substituting
\(u=v-\tau\).  The equality supplies a useful independent check on computed
coefficient blocks.

\subsection{Potential coefficient questions}

The rational functions \(A_m\), or equivalently the integer polynomials
\(P_m\) in \eqref{eq:Am-Pm}, invite further study.  Among the natural questions
are:
\begin{itemize}
\item determining sharp bounds and the zero geometry of \(P_m\);
\item finding direct combinatorial interpretations of its coefficients;
\item deriving a bivariate generating function for the diagonal
      \(R_{m+1,m}\) of \eqref{eq:Rnk-recurrence};
\item locating the exact singularities of \(\Phi_v^{-1}\) and hence the exact
      convergence radius of \eqref{eq:tau-c-series};
\item formalizing the filtered reversion theorem and residual certificates in
      a proof assistant, with coefficient calculations reflected into exact
      rational identities.
\end{itemize}
These questions are not needed for the asymptotic expansion, but they may lead
to a useful polynomial theory associated with the \(r\)-Lambert function.

\section{Conclusions}
\label{sec:conclusion}

The inverse of \(x+\W(x)\) is simple enough to admit an exact generalized
special-function description and rich enough to expose a genuine transseries
phenomenon.

First, the substitution \(u=\W(x)\) turns the inverse problem into
\(u\e^u+u=z\), so \(g(z)=z-\rW_1(z)\).  Second, anchoring at the ordinary
Lambert value \(v=\W(z)\) produces the all-orders expansion
\[
 g(z)=z-v+\sum_{m\ge1}z^{-m}A_m(v),
\]
with the explicit operator formula \eqref{eq:Am-operator}.  Third, the exact
small-parameter equation
\[
 \frac{v}{z}=\frac{v}{v-\tau}-\e^{-\tau}
\]
shows that this outer series converges for sufficiently large \(z\), gives the
remainder scale \((v/z)^{N+1}\), and reveals harmonic-number structure after a
second expansion in \(1/v\).  Finally, the classical inverse-logarithm series
is recovered by expanding \(v\), but that rank-one series necessarily misses
\(\tau\sim(\log z)/z\), which is beyond all its orders.

The general lesson is operational.  For a logarithmic near-identity map
\(F=X+A\), do not guess the inverse coefficient by coefficient in a collection
of unrelated logarithmic monomials.  Place the slow coefficients in a
differential algebra, use the outer variable \(t=X^{-1}\), exploit the fact
that differentiation raises outer valuation, and apply the locally finite
formula
\[
 F^{-1}=X+\sum_{n\ge1}\frac{(-1)^n}{n!}D^{n-1}A^n.
\]
The block formula \eqref{eq:block-formula}, triangular residual cancellation,
Newton iteration, and residual-to-error transfer then provide a complete chain
from formal derivation to certified numerical approximation.

\appendix

\section{Coefficient tables}
\label{app:coefficients}

For reference, write
\begin{equation}
 A_m(v)=\frac{v^{m+1}P_m(v)}{m!(1+v)^{2m-1}}.
 \label{eq:appendix-Pm}
\end{equation}
The first six numerator polynomials are
\begin{align}
 P_1(v)={}&1,\notag\\
 P_2(v)={}&v^2-2,\notag\\
 P_3(v)={}&(2v-1)(v^3-6v-6),\notag\\
 P_4(v)={}&6v^6-8v^5-69v^4-40v^3+96v^2+72v-24,\notag\\
 P_5(v)={}&24v^8-58v^7-428v^6-151v^5+1320v^4+1380v^3\notag\\
 &{}-480v^2-720v+120,\notag\\
 P_6(v)={}&120v^{10}-444v^9-2920v^8+424v^7+16577v^6+17892v^5\notag\\
 &{}-13530v^4-26400v^3-1800v^2+7200v-720.
 \label{eq:Pm-list}
\end{align}
They satisfy
\[
 \deg P_m=2m-2,
 \qquad \operatorname{lc}(P_m)=(m-1)!,
 \qquad P_m(0)=(-1)^{m+1}m!.
\]

The first eight slow Lambert polynomials in
\(\W(z)=L-\ell+\sum_{n\ge1}P_n^{\mathrm{L}}(\ell)L^{-n}\) are
\begin{align}
 P_1^{\mathrm{L}}(y)={}&y,\notag\\
 P_2^{\mathrm{L}}(y)={}&\frac{y(y-2)}2,\notag\\
 P_3^{\mathrm{L}}(y)={}&\frac{y(2y^2-9y+6)}6,\notag\\
 P_4^{\mathrm{L}}(y)={}&\frac{y(3y^3-22y^2+36y-12)}{12},\notag\\
 P_5^{\mathrm{L}}(y)={}&\frac{y(12y^4-125y^3+350y^2-300y+60)}{60},\notag\\
 P_6^{\mathrm{L}}(y)={}&\frac{y(20y^5-274y^4+1125y^3-1700y^2+900y-120)}{120},\notag\\
 P_7^{\mathrm{L}}(y)={}&\frac{y}{840}
 \bigl(120y^6-2058y^5+11368y^4-25725y^3\notag\\
 &\hspace{34mm}{}+24500y^2-8820y+840\bigr),\notag\\
 P_8^{\mathrm{L}}(y)={}&\frac{y}{2520}
 \bigl(315y^7-6534y^6+45962y^5-142149y^4\notag\\
 &\hspace{31mm}{}+205800y^3-135240y^2+35280y-2520\bigr).
 \label{eq:Lambert-polys-8}
\end{align}
The superscript \(\mathrm{L}\) is used only in this appendix to distinguish the
Lambert asymptotic polynomials from the numerator polynomials \(P_m\) in
\eqref{eq:appendix-Pm}.

\section{Derivation of the harmonic correction}
\label{app:harmonic}

This appendix records the coefficient extraction behind
\eqref{eq:cn-harmonic}.  With \(E=1-q\) and
\(T_0=-\log E\), equations \eqref{eq:T1}--\eqref{eq:T2} become
\begin{align}
 T_1={}&-\frac{T_0}{E},\notag\\
 T_2={}&\frac{T_0+T_0^2/2}{E^2}-\frac{T_0^2}{E}.
 \label{eq:T2-T0}
\end{align}
The standard coefficient identities are
\begin{align}
 [q^n]T_0={}&\frac1n,\notag\\
 [q^n]\frac{T_0}{E}={}&H_n,\notag\\
 [q^n]\frac{T_0^2}{E}={}&H_n^2-H_n^{(2)}.
 \label{eq:standard-harmonic-coeffs}
\end{align}
To handle \(T_0^2/E^2\), differentiate:
\begin{equation}
 \frac{\dd}{\dd q}\left(\frac{T_0^2}{E}\right)
 =\frac{2T_0+T_0^2}{E^2}.
 \label{eq:harmonic-derivative}
\end{equation}
Combining \eqref{eq:T2-T0} and \eqref{eq:harmonic-derivative} gives
\begin{equation}
 [q^n]T_2
 =H_n+\frac{n-1}{2}\bigl(H_n^2-H_n^{(2)}\bigr),
 \label{eq:T2-coeff}
\end{equation}
which proves the coefficient of \(v^{-2}\) in \eqref{eq:cn-harmonic}.

\section{Reproducible symbolic computation}
\label{app:code}

The following short SymPy program generates the exact rational functions
\(A_m(v)\) directly from the operator formula \eqref{eq:Am-operator}.  It uses
exact integer and rational arithmetic.

\begin{lstlisting}
import sympy as sp

v = sp.symbols("v")

def theta(expr):
    """The slow derivative v/(1+v) * d/dv."""
    return sp.cancel(v * sp.diff(expr, v) / (1 + v))

def inverse_block(m: int):
    """Return A_m(v) in g=z-W(z)+sum A_m(W(z))/z**m."""
    if m < 1:
        raise ValueError("m must be positive")
    expr = v ** (m + 1)
    for j in range(m):
        expr = sp.cancel(theta(expr) - j * expr)
    return sp.factor((-1) ** (m + 1) * expr / sp.factorial(m + 1))

for m in range(1, 7):
    print(m, inverse_block(m))
\end{lstlisting}

A numerical check may solve the monotone equation \(u\e^u+u=z\) and compare
against successive blocks:

\begin{lstlisting}
import mpmath as mp

mp.mp.dps = 80

def exact_inverse(z):
    z = mp.mpf(z)
    v = mp.lambertw(z)
    u = mp.findroot(lambda y: y * mp.exp(y) + y - z,
                    v - v / z)
    return z - u

# Convert each symbolic A_m to an mpmath callable once.
A = {m: sp.lambdify(v, inverse_block(m), "mpmath")
     for m in range(1, 7)}

def asymptotic_inverse(z, order):
    z = mp.mpf(z)
    vv = mp.lambertw(z)
    result = z - vv
    for m in range(1, order + 1):
        result += A[m](vv) / z**m
    return result

for z in [10, 100, 10_000, 10**8]:
    exact = exact_inverse(z)
    for order in [0, 1, 2, 4, 6]:
        approx = (mp.mpf(z) - mp.lambertw(z)
                  if order == 0
                  else asymptotic_inverse(z, order))
        print(z, order, abs(approx - exact))
\end{lstlisting}

For a production implementation, one should cache polynomial derivatives,
truncate after every arithmetic operation, and verify the final result by the
residual bound \eqref{eq:special-residual-certificate}.

\section{A compact formula sheet}
\label{app:formula-sheet}

For convenience, the principal identities are collected here.

\begin{align*}
&\textbf{Exact reduction:}
&&u=\rW_1(z),\quad u\e^u+u=z,\quad g=z-u.\\[0.4em]
&\textbf{Anchor:}
&&v=\W(z),\quad z=v\e^v,\quad q=\e^{-v}=v/z.\\[0.4em]
&\textbf{Small displacement:}
&&\tau=v-u,\quad q=\dfrac{v}{v-\tau}-\e^{-\tau},\quad g=z-v+\tau.\\[0.4em]
&\textbf{Outer coefficient:}
&&A_m(v)=\dfrac{(-1)^{m+1}}{(m+1)!}
  \prod_{j=0}^{m-1}(\vartheta-j)v^{m+1},
  \quad \vartheta=\dfrac{v}{1+v}\dfrac{\dd}{\dd v}.\\[0.4em]
&\textbf{Outer expansion:}
&&g=z-v+\displaystyle\sum_{m\ge1}\dfrac{A_m(v)}{z^m}.\\[0.4em]
&\textbf{Direct inverse coefficient:}
&&c_n(v)=\dfrac1n[\tau^{n-1}]
  \left(\dfrac{\tau}{v/(v-\tau)-\e^{-\tau}}\right)^n,
  \quad A_n=v^nc_n.\\[0.4em]
&\textbf{Remainder:}
&&g-z+v-\displaystyle\sum_{m=1}^{N}\dfrac{A_m(v)}{z^m}
  =O_N\!\left((v/z)^{N+1}\right).\\[0.4em]
&\textbf{Slow expansion:}
&&g\sim z-L+\ell-\displaystyle\sum_{n\ge1}\dfrac{P_n(\ell)}{L^n}.\\[0.4em]
&\textbf{General reversion:}
&&\bigl(X+A(X)\bigr)^{-1}
  =X+\displaystyle\sum_{n\ge1}\dfrac{(-1)^n}{n!}D^{n-1}A(X)^n.\\[0.4em]
&\textbf{Certification:}
&&\dfrac12|H+\W(H)-z|\le|H-g(z)|\le|H+\W(H)-z|.
\end{align*}

\section*{Acknowledgment of source context}
\addcontentsline{toc}{section}{Acknowledgment of source context}

This note was prepared for the
\emph{series-and-transseries} research-draft area of the ProveIt repository.
The surrounding corpus contains several independent treatments of
polynomial-logarithmic transseries, composition, and series reversal.  The
present article is standalone: all notation, proofs, coefficient formulas, and
verification procedures needed for the specific inverse problem are included
here.

\begin{thebibliography}{99}

\bibitem{CorlessEtAl1996}
R.~M. Corless, G.~H. Gonnet, D.~E.~G. Hare, D.~J. Jeffrey, and D.~E. Knuth,
\emph{On the Lambert \(W\) function},
Advances in Computational Mathematics \textbf{5} (1996), 329--359.
\newblock DOI: \url{https://doi.org/10.1007/BF02124750}.

\bibitem{MezoBaricz2017}
I.~Mez\H{o} and A.~Baricz,
\emph{On the generalization of the Lambert \(W\) function with applications in
theoretical physics},
Transactions of the American Mathematical Society \textbf{369} (2017),
7917--7934.
\newblock DOI: \url{https://doi.org/10.1090/tran/6911}; preprint
\url{https://arxiv.org/abs/1408.3999}.

\bibitem{DLMF413}
NIST Digital Library of Mathematical Functions,
\emph{Section 4.13: Lambert \(W\)-Function},
\url{https://dlmf.nist.gov/4.13}.

\bibitem{Comtet1974}
L.~Comtet,
\emph{Advanced Combinatorics: The Art of Finite and Infinite Expansions},
D. Reidel, Dordrecht, 1974.

\bibitem{FlajoletSedgewick2009}
P.~Flajolet and R.~Sedgewick,
\emph{Analytic Combinatorics},
Cambridge University Press, 2009.

\bibitem{Henrici1974}
P.~Henrici,
\emph{Applied and Computational Complex Analysis}, Vol.~1,
Wiley, New York, 1974.

\bibitem{Olver1997}
F.~W.~J. Olver,
\emph{Asymptotics and Special Functions},
A K Peters, Wellesley, 1997.

\bibitem{Hardy1910}
G.~H. Hardy,
\emph{Orders of Infinity},
Cambridge University Press, 1910.

\bibitem{VanDerHoeven2006}
J.~van der Hoeven,
\emph{Transseries and Real Differential Algebra},
Lecture Notes in Mathematics, Vol.~1888, Springer, 2006.

\bibitem{ADH2017}
M.~Aschenbrenner, L.~van den Dries, and J.~van der Hoeven,
\emph{Asymptotic Differential Algebra and Model Theory of Transseries},
Princeton University Press, 2017.

\bibitem{ProveItRepo}
V.~Reshetnikov,
\emph{ProveIt: series and transseries research drafts},
\url{https://github.com/VladimirReshetnikov/ProveIt/tree/main/Analysis/FabiusFunction/docs/semi-formalized-research-frontiers/drafts/series-and-transseries}.

\end{thebibliography}

\end{document}
